\documentclass[10pt,aspectratio=169]{beamer}
\usepackage{xeCJK}
\setCJKmainfont{HaranoAjiGothic-Regular.otf}[BoldFont=HaranoAjiGothic-Bold.otf]
\setCJKsansfont{HaranoAjiGothic-Regular.otf}[BoldFont=HaranoAjiGothic-Bold.otf]
\usetheme[variant=A]{sciencetokyo}

\title{Presentation Title\\プレゼンテーションのタイトル}
\subtitle{Subtitle}
\author{Taro Kagaku}
\institute{Department of Computer Science, Institute of Science Tokyo}
\date{Event Name, November 23, 2026}

\begin{document}

\maketitle

\begin{frame}{Outline}
  \tableofcontents
\end{frame}

\section{Section I}

\begin{frame}{First Slide}
  \begin{itemize}
  \item Ookayama
    \begin{itemize}
    \item Suzukakedai
      \begin{itemize}
      \item Tamachi
      \end{itemize}
    \end{itemize}
  \item Yushima
  \end{itemize}

  \begin{alertblock}{This is an alert block.}
    This is the content of the block.
  \end{alertblock}

  \begin{exampleblock}{This is an example block.}
    This is the content of the block.
  \end{exampleblock}

  \begin{block}{This is a block.}
    This is the content of the block.
  \end{block}
\end{frame}

\section{Section II}

\begin{frame}{Second Slide: 表と図}
  \begin{columns}
    \begin{column}{0.5\textwidth}
      \begin{table}
        \caption{A Table}
        \begin{tabular}{l|c c}
          \hline
          A & B & C \\
          \hline
          x & y & z \\
          t & u & v \\
          \hline
        \end{tabular}
      \end{table}
    \end{column}

    \begin{column}{0.5\textwidth}
      \begin{figure}
        \begin{tikzpicture}
          \fill[STnavy] (0,0) rectangle (1,2);
          \fill[STlightblue] (1.2,0) rectangle (2.2,1.4);
          \fill[STred] (2.4,0) rectangle (3.4,0.8);
        \end{tikzpicture}
        \caption{A Figure}
      \end{figure}
    \end{column}
  \end{columns}
\end{frame}

\begin{frame}[standout]
  Questions?
\end{frame}

\end{document}
